\documentclass{beamer}
\usepackage{multicol,hyperref}
\renewcommand{\arraystretch}{1.5}
\newcommand{\fttheory}[1]{\fcolorbox{blue}{blue}{\color{white}\textbf{#1}\hspace{\textwidth}}}
\newcommand{\ftexampl}[1]{\fcolorbox{orange}{orange}{\color{black}\textbf{#1}\hspace{\textwidth}}}
\newcommand{\ftdata}[1]{\fcolorbox{violet}{violet}{\color{white}\textbf{#1}\hspace{\textwidth}}}
\newcommand{\ftheader}[1]{\fcolorbox{lightgray}{lightgray}{\color{black}\textbf{#1}\hspace{\textwidth}}}
\begin{document}

\begin{frame}{\ftheader{No-Arbitrage Pricing}}
	\begin{itemize}
		\item \underline{Arbitrage} is the act of costlessly exploiting mispricing of two or more securities to achieve either 1) non-negative profits with certainty or 2) positive profits on average.
		\item Earlier in the course, we studied arbitrage in the context of risk-free bonds.
		\item In this lecture, we are going to study arbitrage in the context of \textit{risky} securities such as stocks, and in so doing, develop the \underline{Arbitrage Pricing Theory} (APT) and lay the groundwork for derivatives pricing
	\end{itemize}
\end{frame}

\begin{frame}
	\begin{center}
		\Huge Alpha
	\end{center}
\end{frame}

\begin{frame}{\fttheory{Alpha}}
	\begin{itemize}
	\item The difference between the stock's actual risk premium and the risk premium of a portfolio with the same systematic risk is the alpha:  
	\begin{equation}
		\alpha_{\text{Stock}} = \underbrace{(\text{E}[R_{\text{Stock}}]-R_{F})}_{\text{True Value}}-\underbrace{\beta_{\text{Stock}}\times(\text{E}[R_{\text{Market}}] - R_{F})}_{\text{CAPM}}
	\end{equation}
	\item While one might observe a non-zero alpha in the short-term, it will likely become zero in the long-term
	\item There are two ways of thinking about $\alpha$, one akin to ``statistical analysis,'' the other to ``fundamental analysis.''
	\end{itemize}
\end{frame}

\begin{frame}{\fttheory{Alpha: Statistical}}
	\begin{itemize}
		\item Run the regression
			\begin{equation}
				r_{S,i}-r_{F,i}=\alpha_{S}+\beta_{S}(r_{M,i}-r_{F,i})+\epsilon_{S,i}
			\end{equation}
			where each $i$ is an observation
		\item If $\alpha_{S}$ is not statistically significant (a $t$-statistic of between $-3$ and $3$), then you (probably) do not have alpha. If it is significant, then either (1) the stock is overpriced (and you should short), or underpriced (and you should long). 
		\item It could also be that your model is wrong. The true model is a multifactor model, and $\alpha_{S}$ is really just the SML (small vs. large) or HML (value vs. growth) factors. 
	\end{itemize}
\end{frame}

\begin{frame}{\fttheory{Alpha: Fundamental}}
	\begin{itemize}
		\item Recall 
			\begin{equation}
				r_{S}-r_{F}=\alpha_{S}+\beta_{S}(r_{M}-r_{F})+\epsilon_{S}
			\end{equation}
	\end{itemize}
\end{frame}

\begin{frame}
	\begin{center}
		\Huge Arbitrage Pricing Theory
	\end{center}
\end{frame}

\begin{frame}{\fttheory{Arbitrage Pricing Theory: Long}}
	Consider a stock that returns
	\begin{equation}
		r_{S}=r_{F}+\alpha+\beta(r_{M}-r_{F})+\epsilon
	\end{equation}
	Importantly, suppose that $\alpha>0$ (the stock is \underline{under}valued). 
	\begin{center}
		\begin{tabular}{l|c|c}
			Security & Cash Flow at $t=0$ & Cash Flow at $t=1$ \\ \hline
			Stock & $-1$ & $(1+r_{S})\times1$ \\
			Market Portfolio & $\beta$ & $-(1+r_{M})\times\beta$ \\
			Risk-Free Security & $1-\beta$ & $-(1+r_{F})\times(1-\beta)$\\ \hline
			Arbitrage Portfolio & $0$ & $\alpha+\epsilon$\\
		\end{tabular}
	\end{center}
	Here, we have a portfolio that pays $E[\alpha+\epsilon]=\alpha>0$ on average and costs nothing to setup. 
\end{frame}

\begin{frame}{\ftexampl{Example: Long}}
	Walmart has a beta of $.3659$ and an alpha of $.5450\%$.
	\begin{center}
		\begin{tabular}{ll}
			\multicolumn{2}{l}{Date $t=0$} \\ \hline
			1. & Short .3659 USD worth of market.\\
			2. & Borrow .6341 USD.\\
			3. & Use your 1 USD to buy 1 USD worth of Walmart. \\ \hline
			\multicolumn{2}{l}{Date $t=1$: Suppose $r_{M}=.5\%$, $r_{S}=1\%$ and $r_{F}=.1\%$} \\ \hline
			1. & Sell Walmart and receive $(1+.01)\times 1=1.01$ USD. \\
			2. & Repay the loan, which costs $(1+.001)\times.6341=.6347$ USD \\
			3. & Cover your short, which costs $(1+.005)\times.3659=.3677$ \\ \hline
			\multicolumn{2}{l}{You earned .0076 USD. You should earn .005450 USD on \underline{average}.}
		\end{tabular}
	\end{center}
\end{frame}

\begin{frame}{\fttheory{Arbitrage Pricing Theory: Short}}
	Same setup as before...
	\begin{equation}
		r_{S}=r_{F}+\alpha+\beta(r_{M}-r_{F})+\epsilon
	\end{equation}
	...but now suppose that $\alpha<0$ (the stock is \underline{over}valued). 
	\begin{center}
		\begin{tabular}{l|c|c}
			Security & Cash Flow at $t=0$ & Cash Flow at $t=1$ \\ \hline
			Stock & $1$ & $-(1+r_{S})\times 1$ \\
			Market Portfolio & $-\beta$ & $(1+r_{M})\times\beta$ \\
			Risk-Free Security & $\beta-1$ & $-(1+r_{F})\times(\beta-1)$\\ \hline
			Arbitrage Portfolio & $0$ & $-\alpha+\epsilon$\\
		\end{tabular}
	\end{center}
	We have a portfolio that pays $E[-\alpha+\epsilon]=-\alpha>0$ on average and costs nothing to setup. 
\end{frame}

\begin{frame}{\ftexampl{Example: Short}}
	Walgreen's has a beta of $1.0733$ and an alpha of $-.3471\%$.
	\begin{center}
		\begin{tabular}{ll}
			\multicolumn{2}{l}{Date $t=0$} \\ \hline
			1. & Short 1 USD worth of Walgreen's. \\
			2. & Borrow .0733 USD. \\
			3. & Use your 1.0733 USD to buy 1.0733 USD worth of market. \\ \hline
			\multicolumn{2}{l}{Date $t=1$: Suppose $r_{M}=1\%$, $r_{S}=.5\%$ and $r_{F}=.1\%$} \\ \hline
			1. & Sell the market and receive $(1+.01)\times 1.0733=1.084$ USD. \\
			2. & Repay the loan, which costs $(1+.001)\times .0733=.0734$ USD. \\
			3. & Cover your short, which costs $(1+.005)\times 1 = 1.005$ USD. \\ \hline
			\multicolumn{2}{l}{You earned .0056 USD. You should earn .003471 USD on \underline{average}.}
		\end{tabular}
	\end{center}
\end{frame}

\begin{frame}{\fttheory{The Arbitrage Pricing Theory}}
	\begin{itemize}
		\item These examples show that if $\alpha\neq0$, then one can costlessly construct a portfolio that pays $|\alpha|$ on average. 
		\item Suppose, for concreteness, that $\alpha>0$ (as in the Walmart example). Then arbitrageurs start buying the stock. 
		\item According to the laws of supply and demand, the price will rise (the return will fall) until $\alpha=0$.
		\item The APT claims that $\alpha=0$ on average:
			\begin{equation}
				\text{E}[r_{S}-r_{F}]=\beta\text{E}[r_{M}-r_{F}]
			\end{equation}
			which is exactly what the CAPM claims. 
		\item The APT tells us \textit{why} the CAPM should hold and \textit{how} to profit when it breaks.
		\item Note that these trades involve idiosyncratic risk ($\epsilon$) but \underline{not} systematic risk ($r_{M}$)
	\end{itemize}
\end{frame}

\begin{frame}{\ftdata{Long-Short Equity}}
	\begin{itemize}
		\item Hedge funds that employ trades like the one in the Walmart example are called \underline{long-short equity} funds
		\item They are often called \underline{market-neutral} funds, because their returns aren't correlated with those of the market
		\item Some long-short equity funds have a ``long-bias'' and follow a 130/30 rule: they are 130\% long and 30\% short (put differently, 30\% of their long are financed with shorts)
		\item Although we have considered trades in which the arbitrageur goes long/short the stock and short/long the market, many long-short equity funds will employ \underline{paired trades} in which they go long one stock (say Coca-Cola) and short another (say Pepsi) with the belief that the long stock with outperform the short stock, even if both stocks rise or fall
	\end{itemize}
\end{frame}

\end{document}
